\documentclass[letter,11pt]{article} % Tamaño de página y letra, tipo de documento
\usepackage[top = 2.0cm, bottom = 2.0cm, left = 2.0cm, right = 2.0cm]{geometry}

%% Comandos para codificación de archivo
\usepackage[utf8]{inputenc}
\usepackage[spanish,es-tabla]{babel}
\spanishdecimal{.}

%%%%%%%%%%%%%%%%%%%%%%%%%%%%%%%%%%%%%%%%%%%%%%%%
%% Bibliotecas útiles %%
\usepackage{multirow} % Múltiples renglones
\usepackage{multicol} % Múltiples columnas
\usepackage{booktabs} % Tablas más estéticas
\usepackage{graphicx} % Agregar figuras
\usepackage{setspace}  % Quita sangría de inicio de párrafos
\setlength{\parindent}{0in}
\usepackage{float} % Control de ubicación de figuras
\usepackage{fancyhdr} % Formato de encabezados
\usepackage{listings} % Para código fuente
\usepackage{amsmath,mathtools,amsfonts,amssymb} % Símbolos matemáticos
\usepackage{xcolor} % Texto y ecuaciones en color
\usepackage{pdfpages} % Insertar PDFs
\usepackage{hyperref}
\newcommand{\Heaviside}{\mathrm{H}}
%%%%%%%%%%%%%%%%%%%%%%%%%%%%%%%%%%%%%%%%%%%%%%%%
%% Encabezado y pie de página %%
\pagestyle{fancy}  
\fancyhf{} 
%%%%%%%%%%%%%%%%%%%%%%%%%%%%%%%%%%%%%%%%%%%%%%%%
\lhead{\footnotesize Simulación Rover Tarea 04} 
\rhead{\footnotesize LIRA} 
\cfoot{\footnotesize \thepage} 
%%%%%%%%%%%%%%%%%%%%%%%%%%%%%%%%%%%%%%%%%%%%%%%%
%% Aquí comienza la edición del documento

\begin{document}
	\thispagestyle{empty} % Sin encabezado en la primera página
	
	%%%%%%%%%%%%%%%%%%%%%%%%%%%%%%%%%%%%%%%%%%%%%%%%
	%% Datos para la primera página
	\begin{tabular} {p{15cm} p{2cm}}
		&  \multirow{5}{*}{\includegraphics[scale=0.28]{UNAM_logo.png}}\\
		{\large \bf Laboratorio de Innovación y Robótica Avanzada} & \\
		{\large \bf LIRA} & \\
		UNAM | FI & \\
		&  \\
		\hline
	\end{tabular} 
	
	\vspace*{0.3cm} 
	
	\begin{center}
		{\Large \bf Tarea 04} \\
		{\bf {\normalsize Reconocimiento de objetos mediante visión parte 1}} \\
		\vspace{2mm} 
	\end{center}
	
	%Documento
	
	\section{Visión artificial}
	
	\textbf{¿Que es la visión artificial?}
	
    En la inteligencia artificial la visión es la parte que entrena a computadoras y sistemas para identificar, procesar e interpretar datos visuales del mundo real mediante imágenes y vídeos, de manera similar a la visión humana. Utiliza algoritmos de machine learning y redes neuronales para analizar patrones y tomar decisiones automatizadas con alta precisión. \\

     \\
    
	\textbf{Ejemplos de entrenamiento para el rover:}
    \begin{enumerate}
        \item {Evasión de obstaculos como rocas}
        \item {Detección y clasificación de objetos}
        \item {Navegación autónoma}
        \item {Reconocimiento de terrenos}
        \item {Seguimiento de objetivos}
        \item {Inspección visual}
    \end{enumerate}
	
	\section{Problema para el rover}
	
	El objetivo de esta tarea es aprender a utilizar una computadora con ayuda de una cámara para identificar objetos que se puedan presentar en nuestro camino, el ejemplo que utilizaremos sera una roca como la de la imagen.\\ 

    \begin{figure}[H]
			\centering
			\includegraphics[width=0.2\linewidth]{Piedra.jpeg}
            \caption{Ejemplo de piedra}
			\label{seq}
		\end{figure}

    \section{Instalación de Labelme}

    Para evitar conflictos se debe untilizar un entorno virtual 
    
    \begin{enumerate}
        \item {(Ubuntu 24.04) Abrir una terminal}
        \item {Instalación del entorno virtual}
        \begin{lstlisting}
           sudo apt install python3-venv 
        \end{lstlisting} 
        \item {Crear y activar el entorno virtual}
        \begin{lstlisting}
           python3 -m venv <your-venv-name>        #crea venv
           source <your-venv-name>/bin/activate    #activa venv
        \end{lstlisting} 
        \item {Dentro del entorno virtual se instala LabelMe}
        \begin{lstlisting}
           pip install labelme 
        \end{lstlisting} 
        \item {Ejecutar aplicación}
        \begin{lstlisting}
           labelme 
        \end{lstlisting} 
    \end{enumerate}
    
    \section{Selección de objetos para entrenamiento}
        \begin{enumerate}
            \item Abrir el directorio que contiene las imagenes
            (PAPIME\_PE103825/Ejercicios/Ejercicio04/Im\_Entrenamiento)
            \begin{figure}[H]
    			\centering
    			\includegraphics[width=0.6\linewidth]{Directorio.png}
                \caption{Selección carpeta}
    			\label{seq}
    		\end{figure}
                \begin{figure}[H]
        			\centering
        			\includegraphics[width=0.6\linewidth]{Carpeta.png}
                    \caption{Selección carpeta}
        			\label{seq}
        		\end{figure}
            \item Seleccionar la herramienta Create Rectangle ubicada en el menú de herramientas del lado izquierdo de la pantalla
            \item Encerrar en un rectángulo el objeto a reconocer lo más justo posible.
            \item Nombrar el objeto y seleccionar ok.
            \begin{figure}[H]
    			\centering
    			\includegraphics[width=0.5\linewidth]{Instrucciones.png}
                \caption{Nombre de objeto}
    			\label{seq}
    		\end{figure}
            \item Seleccionar Next Image del menú de herramientas superior
            \item Selecciona save y cuando aparezca el menú de la imagen hacer clck en save nuevamente
            \begin{figure}[H]
        			\centering
        			\includegraphics[width=0.6\linewidth]{Guardar.png}
                    \caption{Guardar}
        			\label{seq}
        		\end{figure}
            \item Repetir este proceso con todas las imagenes de la carpeta
            \item Mover todos los archivos .json a una carpeta nueva que solo contenga estos archivos
            \\
            \\
            cd PAPIME\_PE103825/Ejercicios/Ejercicio04/Im\_Entrenamiento
            \\
            mv *.json /home/\<usuario\>/PAPIME\_PE103825/Ejercicios/Ejercicio04/Im\_Entrenado
        \end{enumerate}

        \section{Funciones básicas de Labelme}
        \textbf{Edit: }Crear figuras para las etiquetas, Ctrl + R para crear rectángulos. \\

        \textbf{Next Image: }Hacer clic después de guardar los cambios. Si no se guardan, el sistema lo solicitará.\\
        
        \textbf{Save: }Guarda los cambios y crea un archivo .json con el mismo nombre de la imagen, no lo cambies.
        Debe existir una imagen con su archivo .json correspondiente.\\
        
        \textbf{Edit Polygons: }Si alguna forma o etiqueta es incorrecta, puedes corregirla con esta herramienta.\\
        
        \textbf{Label List: }En esta sección puedes verificar todas las etiquetas que has creado. Es muy importante que cada etiqueta coincida correctamente con su polígono correspondiente. Si hay errores o desajustes, el modelo no aprenderá de manera efectiva.\\
        
        \textbf{Polygon Labels: }Durante la sesión actual de LabelMe, esta sección muestra todas las etiquetas que has asignado a las imágenes. Algunas imágenes pueden contener menos objetos para etiquetar que otras, lo cual se reflejará aquí. Las etiquetas marcadas indican las que se utilizan en la imagen actual.\\
        
        \textbf{File List: }En esta sección puedes ver cuántas imágenes ya han sido etiquetadas y cuáles no. Las casillas marcadas indican que la imagen tiene etiquetas, mientras que las casillas sin marcar significan que la imagen aún no tiene ninguna etiqueta.
	
	
    
    \section{Experimentación}
		Una vez observado el funcionamiento del programa labelme y haber etiqutado las fotos de la carpeta puedes conseguir tu propia piedra, tomarle fotos y etiqutarlas para que funcione mejor el reconocimiento, mientrtas más fotos y posiciones se coloquen será mejor el entramiento que se hará posteriormente.
	
\end{document}